\documentclass[11pt,a4paper]{article}
\usepackage{amsmath,amssymb,amsthm}
\usepackage[margin=2.5cm]{geometry}
\usepackage{hyperref}

\newtheorem{definition}{Definition}[section]
\newtheorem{theorem}[definition]{Theorem}
\newtheorem{proposition}[definition]{Proposition}
\newtheorem{lemma}[definition]{Lemma}
\newtheorem{corollary}[definition]{Corollary}
\newtheorem{conjecture}[definition]{Conjecture}
\newtheorem{remark}[definition]{Remark}

\title{Hyperseed Formalization:\\
  Proof of Humanity: Beyond the Orb}
\author{ProtomegaTron (automated formalization)}
\date{2026-09-17}

\begin{document}
\maketitle

\begin{abstract}
We formalize the intellectual structure of Ben Goertzel's Substack article
``Proof of Humanity: Beyond the Orb'' (2026-08-26) within the Hyperseed
ontology. The article proposes an open, multi-vendor, role-based
proof-of-humanity protocol as an alternative to Worldcoin's single-device
approach. Each architectural principle maps onto Hyperseed structures:
composite proof as sheaf of evidence sections, role-specific certification
as fiber decomposition of trust, the protocol-not-product design as
distributed fiber bundle, media provenance parallels as $E_\pi$ applied to
identity, three reputation forms as three-tier holonomy, and anti-monoculture
as anti-trivial-cocycle.
\end{abstract}

\tableofcontents

%% =================================================================
\section{Monoculture as trivial cocycle}
\label{sec:monoculture}
%% =================================================================

\begin{definition}[Single-vendor monoculture --- claims 1--4, 41]
\label{def:monoculture}
Let $\mathcal{D}$ be the space of proof-of-humanity devices. The Worldcoin/World
approach restricts $\mathcal{D}$ to a single device $d_{\text{Orb}}$:
\[
  \mathcal{D}_{\text{World}} = \{d_{\text{Orb}}\}.
\]
All transition functions between evidence sources become trivial:
$g_{ij} = \text{id}$ for all $i, j$, because there is only one source to
compare with itself. This is a \emph{trivial cocycle} --- the system cannot
detect its own failure because there is no independent perspective from
which to triangulate.
\end{definition}

\begin{proposition}[Correlated failure --- claim 4]
\label{prop:correlated}
In a trivial-cocycle system, failure is maximally correlated: if
$d_{\text{Orb}}$'s biometric is broken (by AI synthesis, hardware flaw,
or supply-chain attack), the entire network falls simultaneously. There
is no redundancy at the evidence-source level, only at the instance level
(many Orbs, all with the same vulnerability).
\end{proposition}

\begin{corollary}[Vendor lock-in as oligarch capture]
The single-vendor architecture creates the same structural capture that
the decentralized community exists to resist --- the identity layer of
the entire ecosystem depends on one corporation's hardware decisions,
security practices, and business continuity.
\end{corollary}

%% =================================================================
\section{Composite proof as sheaf of evidence sections}
\label{sec:composite}
%% =================================================================

\begin{definition}[Evidence sheaf --- claims 5--7, 36]
\label{def:evidence-sheaf}
A \emph{proof of humanity} is a global section of an evidence sheaf
$\mathcal{E}$ over the space of evidence roles $R$:
\begin{itemize}
\item Each role $r_k \in R$ (iris acquisition, presentation-attack detection,
  biological presence, secure display, anti-relay, action authorization)
  defines an open set $U_k \subset R$,
\item Each certified device $d_i$ for role $r_k$ produces a local section
  $e_{i,k} : U_k \to \text{Evidence}$,
\item The global proof is assembled via descent data: transition functions
  $g_{ij}$ that verify consistency between different devices' evidence
  for overlapping roles.
\end{itemize}
``Human'' is not a single sensor output --- it is a policy-selected
composition of narrowly scoped claims.
\end{definition}

\begin{proposition}[Same as media provenance --- claim 6]
\label{prop:media}
Validating that a human was present in a certain place doing a certain
thing is the same sheaf-assembly problem as media provenance: composite
proof from camera hardware, GPS, secure elements, surrounding context,
adjudicated by different parties. Proof of humanity is $E_\pi$ restricted
to the identity fiber (claim 39).
\end{proposition}

\begin{theorem}[Nontrivial cocycle advantage]
\label{thm:nontrivial}
A multi-device, multi-vendor evidence sheaf with nontrivial transition
functions $g_{ij} \neq \text{id}$ has strictly greater failure detection
capability than a trivial-cocycle monoculture:
\begin{enumerate}
\item \textbf{Cross-validation}: inconsistency between devices $d_i$ and
  $d_j$ detecting the same property from different physical principles
  reveals compromise of either device.
\item \textbf{Graceful degradation}: revocation of one device replaces
  one local section while preserving the global section (claim 43).
\item \textbf{Evolutionary adaptability}: a new sensor technology enters
  as a new role profile without invalidating existing sections (claim 11).
\end{enumerate}
\end{theorem}

%% =================================================================
\section{Role-specific certification as fiber decomposition of trust}
\label{sec:certification}
%% =================================================================

\begin{definition}[Trust fiber --- claims 9, 16--18, 37]
\label{def:trust-fiber}
The trust bundle $\pi_T : E_T \to \mathcal{D}$ over the space of devices
has a fiber that decomposes into independent role components:
\[
  E_{T,d} = \bigoplus_{k=1}^{|R|} T_{d,k}
\]
where $T_{d,k}$ is device $d$'s certification status for role $r_k$.
Key properties:
\begin{itemize}
\item \textbf{Role-specific}: a device is certified only for the roles
  its testing actually supports, not for ``being trustworthy'' in general.
\item \textbf{Configuration-specific}: certification is bound to exact
  model, hardware revision, firmware measurement, algorithm version,
  calibration state, and operating environment.
\item \textbf{Independent verification}: certification comes from
  accredited independent laboratories, not manufacturer self-certification.
\end{itemize}
This is the same unbundling principle as the AI-rights fiber decomposition
(welfare/identity/standing/participation/franchise).
\end{definition}

\begin{proposition}[Section replacement --- claims 18, 43]
\label{prop:replacement}
When a device $d_i$ is compromised, its certification for role $r_k$ is
revoked: $T_{d_i, k} \mapsto \emptyset$. The global proof is repaired by
substituting an alternative certified device $d_j$ for the same role:
$e_{i,k} \mapsto e_{j,k}$. The person's identity persists because identity
is tied to the global section, not to any particular local section.
\end{proposition}

%% =================================================================
\section{Protocol not product as distributed fiber bundle}
\label{sec:protocol}
%% =================================================================

\begin{definition}[Distributed protocol --- claims 8, 10--11, 38]
\label{def:protocol}
The proof-of-humanity protocol is a distributed fiber bundle: no single
device or vendor provides a global section. Instead:
\begin{itemize}
\item Many manufacturers produce devices implementing open capability
  profiles (local sections over different open sets),
\item The verifier selects a policy (a covering of the role space $R$),
\item The protocol provides the descent data (standards for cross-device
  evidence consistency).
\end{itemize}
This is the same distributed fiber bundle architecture as:
\begin{enumerate}
\item Decentralized AGI development (How Omega Lost Its Claw),
\item The positive left AGI agenda (Sanders article),
\item Open-source AGI support (AI Rights article).
\end{enumerate}
\end{definition}

\begin{remark}[Policy as covering choice]
Different assurance levels correspond to different coverings of $R$:
a high-assurance policy requires many overlapping open sets (multiple
devices, multiple roles, cross-device binding), while a low-stakes policy
requires minimal covering (one uniqueness credential). The protocol is
invariant; only the covering changes.
\end{remark}

%% =================================================================
\section{Identity architecture: fiber-local identity}
\label{sec:identity}
%% =================================================================

\begin{definition}[Domain nullifiers as fiber-local identity --- claims 19--22, 42]
\label{def:nullifiers}
Cryptographic domain nullifiers implement identity sections that are
local to their domain:
\begin{itemize}
\item Within domain $D_\alpha$: the nullifier proves ``this is the same
  person who previously participated in $D_\alpha$'' (the sheaf condition
  holds within the domain),
\item Across domains $D_\alpha, D_\beta$: the nullifiers for different
  domains are cryptographically unlinkable (no identity information leaks
  between domains).
\end{itemize}
Civil identity can be selectively layered on where law requires it,
rather than being welded into the proof of humanity itself.
\end{definition}

\begin{proposition}[Pseudonymous reputation as directed history]
\label{prop:pseudonymous}
A pseudonymous participant accumulates reputation as a directed history
within each domain:
\[
  \mathcal{H}_{p, D_\alpha} = (h_1, h_2, \ldots, h_n)
\]
where each $h_i$ records an action, contribution, or evaluation within
domain $D_\alpha$. The directed history provides continuity (``this is the
same scientist who published that paper'') without requiring civil identity
disclosure. This is the same versioned-identity structure from the AI Rights
article, applied to pseudonymous humans.
\end{proposition}

%% =================================================================
\section{Three reputation forms as three-tier holonomy}
\label{sec:reputation}
%% =================================================================

\begin{definition}[Three-tier reputation holonomy --- claims 24--26, 40]
\label{def:three-tier}
The democratic trust fabric comprises three interlinked reputation tiers:
\begin{enumerate}
\item \textbf{Human reputation} $\text{Rep}_H$: the directed history of
  a human participant's contributions, endorsements, and track record.
\item \textbf{Device/certifier reputation} $\text{Rep}_D$: the track
  record of devices and certification laboratories --- how often their
  evidence has been confirmed, challenged, or revoked.
\item \textbf{AI agent reputation} $\text{Rep}_A$: the performance
  history of AI agents acting under human delegation, including what
  they were authorized to do and what they actually did.
\end{enumerate}
These form a three-tier holonomy structure:
\[
  \text{Rep}_H \xrightarrow{\text{certifies}} \text{Rep}_D
  \xrightarrow{\text{grounds}} \text{Rep}_A
  \xrightarrow{\text{serves}} \text{Rep}_H,
\]
where each tier's cocycle data feeds into the next tier's curvature
assessment. This is analogous to the cocycle\_status $\to$ End($E_c$)
holonomy $\to$ curvature hierarchy from the three-tier holonomy note.
\end{definition}

\begin{proposition}[Agent authorization binding --- claim 25]
\label{prop:agent-auth}
High-impact agent actions require a recent composite humanity proof
tied to the exact action:
\[
  \text{auth}(a) = (\text{PoH}_{\text{recent}}, \text{action}_a,
  \text{scope}_a, \text{expiry}_a)
\]
This binds which human credential authorized which agent, under which
scope and expiry, without publishing the human's biometrics or civil
identity. The authorization is a local section of the trust bundle
that connects the human and agent reputation tiers.
\end{proposition}

%% =================================================================
\section{Protocol is not a coin}
\label{sec:not-coin}
%% =================================================================

\begin{proposition}[Modular payment as adapter layer --- claims 27--29]
\label{prop:payment}
Tying the protocol to one cryptocurrency creates the same monoculture
risk at the economic layer that the Orb creates at the hardware layer:
single point of failure, vendor lock-in, correlated failure. Instead:
\begin{itemize}
\item The protocol exposes modular payment and anchoring adapters,
\item Different deployments use different economic mechanisms
  (token, procurement contract, subsidy),
\item All produce the same open evidence objects.
\end{itemize}
``One's status as a human shouldn't fluctuate with one market'' is
the economic analog of ``no single device is indispensable'' ---
the anti-trivial-cocycle principle applied to the payment layer.
\end{proposition}

%% =================================================================
\section{Cross-references}
%% =================================================================

\begin{remark}[Connections to other formalizations]
\begin{itemize}
\item \textbf{AI Rights} (2026-09-02): unbundled rights as fiber
  decomposition. Here the same principle applied to device certification
  and identity properties. The copy problem's cryptographic civic identity
  is the AI-side of the same infrastructure.
\item \textbf{How Omega Lost Its Claw} (2026-09-11): distributed fiber
  bundle as the architecture for decentralized systems. Here applied to
  identity rather than AGI development.
\item \textbf{Coxon/EA} (2026-09-13): auditor capture as cocycle failure.
  Here the certification system is designed to prevent it via independent
  labs and role-specific testing.
\item \textbf{Sanders Ban} (2026-09-05): selection effect as provenance
  destruction. The Orb monoculture is a gentler version: not destroying
  provenance, but making it trivial (only one source to compare).
\item \textbf{Seven Flavors} (2026-07-01): emergent stupidity as failed
  descent data. The Orb's trivial cocycle prevents the multi-perspective
  descent data that would catch systemic failures.
\item \textbf{Three-tier holonomy note} (LTM 568f40b4): cocycle\_status
  $\to$ End($E_c$) $\to$ curvature. Here: human reputation $\to$
  device/certifier reputation $\to$ AI agent reputation.
\end{itemize}
\end{remark}

\begin{conjecture}[Minimum covering for robust proof]
What is the minimum number of independent roles (open sets in $R$)
needed for a proof of humanity to be robust against single-device
compromise? Conjecture: at least three independent physical principles
(e.g., iris biometric, biological presence via a different modality,
and behavioral challenge-response), providing a covering of $R$ with
enough overlap that any single compromised section is detectable by
cross-validation with the other two.
\end{conjecture}

\begin{conjecture}[Reputation convergence]
In the three-tier holonomy, does the reputation of AI agents
($\text{Rep}_A$) converge to reflect the reputation of their authorizing
humans ($\text{Rep}_H$) over time? Conjecture: partially --- agent
reputation accumulates independently but is bounded by the trust
ceiling of the authorizing human, creating a semi-independent
directed history that inherits initial conditions from the human tier
but can diverge.
\end{conjecture}

\end{document}
